\documentclass[10pt,a4paper]{amsart}
\usepackage[margin=19mm]{geometry}
\usepackage{amsmath,amssymb,mathtools}
\usepackage[scr=rsfs,cal=euler]{mathalfa}
\usepackage{xfrac}
\usepackage[unicode,hidelinks,bookmarks=false]{hyperref}
\newcommand{\ie}{i.e.\ }
\newcommand{\cf}{cf.\ }
\newcommand{\resp}{respectively\ }
\newcommand{\Subsection}{\subsection}
\providecommand{\nobreakdash}{\nobreak\hbox{-}}
\setlength{\emergencystretch}{2em}
\multlinegap0pt
\usepackage{bm}
\usepackage{bbm}
\usepackage{dsfont}
\usepackage{xspace}
\usepackage{cancel}
\usepackage{mathtools}
\usepackage{amsthm}
\makeatletter
\@ifundefined{Definition}{\newtheorem{Definition}{Definition}}{}
\@ifundefined{Lemma}{\newtheorem{Lemma}[Definition]{Lemma}}{}
\@ifundefined{Proposition}{\newtheorem{Proposition}[Definition]{Proposition}}{}
\makeatother
\newcommand{\N}{\mathbb N}
\newcommand{\g}{\gamma}
\newcommand{\hps}{\hat{\psi}}
\newcommand{\hvp}{\hat{\varphi}}
\newcommand{\ps}{\psi}
\newcommand{\s}{\sigma}
\newcommand{\vp}{\varphi}
\title[Timelike conjugate points in Lorentzian length spaces]{Timelike conjugate points in Lorentzian length spaces}
\author{James D.E. Grant, Michael Kunzinger, Argam Ohanyan, Yasmin Schinnerl, Roland Steinbauer}
\date{Source: arXiv:2509.12855v3 (CC BY 4.0)}
\begin{document}
\maketitle

\noindent\emph{Source and licence.} Timelike conjugate points in Lorentzian length spaces. Version arXiv:2509.12855v3 (CC BY 4.0), distributed under the Creative Commons Attribution 4.0 International licence, as recorded in the arXiv licence metadata for this exact version and shown on the abstract page https://arxiv.org/abs/2509.12855v3. Full rights notice: licenses/arxiv-2509.12855-CC-BY-4.0.txt.\par\smallskip

\noindent\textit{Editorial scope.} Counted unit: the Proposition whose source label is \texttt{Prop:Frechet-distance} in the article (source0.tex lines 1153--1154, source label \texttt{Prop:Frechet-distance}), delivered together with the article's own complete original proof of it (source0.tex lines 1155--1211, one \texttt{proof} environment of 57 lines). No mathematics is written, repaired or re-derived: statement and proof are the article's own text line for line. Required context retained uncounted: Lem:parametrizations-Hausdorff-compact (lines 1134--1137). Every definition, display and label of those lines is retained unchanged. Omitted from this package and not redistributed: 1--1133, 1138--1152, 1212--1358. Nothing inside a retained excerpt is omitted or altered: every retained body is delivered exactly as the article's lines are written, so no omission receipt of any kind accompanies this package. Citations: the retained excerpts carry no citation command at all, so no bibliography entry of the article is transcribed and no reference is redistributed. Reference adaptation: none was needed and none was made; every reference inside the delivered unit resolves inside it, so no unresolved reference or citation is produced. Printed slips kept literal: no printed word, symbol, hypothesis or number of the counted statement or of its proof was altered. Layout and class: the article's own class is \texttt{article} with packages including inputenc, appendix, a4wide, amssymb, amsmath, amsthm, mathtools, euscript, ulem, hyperref, xcolor, graphicx, changebar, enumitem; the wrapper is the lane's pinned \texttt{amsart} editorial wrapper at 10pt on A4, which restates the theorem-like environments the retained text needs and adds the article's own macro definitions that the retained text needs (\texttt{\textbackslash N}, \texttt{\textbackslash g}, \texttt{\textbackslash hps}, \texttt{\textbackslash hvp}, \texttt{\textbackslash ps}, \texttt{\textbackslash s}, \texttt{\textbackslash vp}), with extra wrapper packages (bm, bbm, dsfont, xspace, cancel, mathtools). The delivered numbering is the unit's own. Assets: no retained span contains a figure environment, a graphics-inclusion command, a PDF-page inclusion, a table or a TikZ picture; no image file is copied into this package, the package declares no asset dependency and no image file is redistributed with it. Licence: version arXiv:2509.12855v3 of this article is distributed under the Creative Commons Attribution 4.0 International licence, as recorded in the arXiv licence metadata for this exact version and shown on the abstract page https://arxiv.org/abs/2509.12855v3. A full rights notice is delivered at licenses/arxiv-2509.12855-CC-BY-4.0.txt. Characters: every retained excerpt is pure ASCII, so the delivered engine, XeLaTeX with its default Latin Modern text font, reports no missing character.
\begin{Lemma}\label{Lem:parametrizations-Hausdorff-compact}
Let $N$ be the set of continuous, coordinate-wise non-decreasing maps $\alpha=(\varphi,\psi) \colon [0,1]\to [0,1]^2$ with $\alpha(0) = (0,0)$
and $\alpha(1)=(1,1)$. Then $\mathcal{N} \coloneqq \{\alpha([0,1]) \mid \alpha\in N\}$ is compact with respect to the Hausdorff distance $d_H$.
\end{Lemma}

\begin{Proposition}\label{Prop:Frechet-distance} The Fr\'echet distance $d_F$ defines a metric on the space of nowhere constant curves in $(X,d)$.  
\end{Proposition}

\begin{proof} Since parametrizations form a group, the Fr\'echet distance is well defined. Symmetry of $d_F$ is clear and the triangle inequality follows readily from the definition. Also $d_F([\gamma],[\gamma])=0$ is obvious. 

The only non-trivial point therefore is to show that $d_F([\gamma],[\sigma])=0$ implies $[\gamma]=[\sigma]$, i.e., that $\gamma$ is a reparametrization of $\sigma$. From the definition of $d_F$ we obtain sequences of parametrizations $(\varphi_n,\psi_n)$ such that, for each $n\in\N$:
\begin{equation}\label{eq:phin-psin}
    \max_{t\in [0,1]} d(\gamma(\varphi_n(t)),\sigma(\psi_n(t))) < \frac{1}{n}.
\end{equation}

Set $\alpha_n \coloneqq (\varphi_n,\psi_n)$, and $A_n \coloneqq \alpha_n([0,1])$. Then by Lemma~\ref{Lem:parametrizations-Hausdorff-compact}, $(A_n)_{n\in \N}$
is a sequence in the $d_H$-compact set $\mathcal{N}$, hence possesses a $d_H$-convergent subsequence, which we may without loss of generality assume to be $(A_n)_{n\in \N}$ itself. Let $A\in \mathcal{N}$ be its limit. Then there is some $\alpha = (\varphi,\psi)\in N$ with $A=\alpha([0,1])$. Given any $t\in [0,1]$, there exists a sequence $(t_n)_{n\in \N}$ in $[0,1]$ such that 
\[
A_n\ni (\varphi_n(t_n),\psi_n(t_n)) \to (\varphi(t),\psi(t)) \qquad (n\to \infty).
\]
By \eqref{eq:phin-psin} we have $d(\gamma(\varphi_n(t_n)),\sigma(\psi_n(t_n))) < \frac{1}{n}$, so letting $n\to \infty$ we obtain
$d(\gamma(\varphi(t)),\sigma(\psi(t)))=0$, implying that $\gamma\circ \varphi = \sigma\circ \psi$. 

We are left with the task of constructing from $\varphi$ and $\psi$ new continuous surjective maps $\hat\varphi$, $\hat\psi \colon [0,1]\to [0,1]$ that are strictly monotonically increasing, hence qualify as parametrizations and satisfy $\gamma\circ \hat\varphi = \sigma\circ \hat\psi$. This we do
in several steps.

We first claim that $\varphi$ is constant on any non-trivial interval $[t_1,t_2]$ of $[0,1]$ if and only if $\psi$ is constant on $[t_1,t_2]$.
By symmetry, it suffices to show the `only if' part of the claim. Since $\varphi$ is constant on $[t_1,t_2]$, so is
$\gamma\circ \varphi = \sigma\circ \psi$.
Also, $\psi$ being continuous and non-decreasing, the image of $[t_1, t_2]$ under $\psi$ is the closed interval $J = [\psi(t_1), \psi(t_2)]$.
The path $\sigma$ is therefore constant on $J$. By our assumption on $\sigma$ this forces $J$ to be a single point, hence
$\psi(t_1) = \psi(t_2)$. Since $\psi$ is non-decreasing, $\psi$ is indeed constant on $[t_1, t_2]$.

Define $\tau \colon [0,1] \to [0,1]$ by
$$
\tau(t) = \frac{\varphi(t) + \psi(t)}{2}.
$$
Then $\tau$ is continuous, non-decreasing and surjective $[0,1]\to [0,1]$. 

Next we show that $\tau(t)$ is constant on an interval $[t_1, t_2]$ if and only if both $\vp(t)$ and $\ps(t)$ are constant on that interval. For the non-trivial direction of this claim, suppose that $\tau$ is constant on $[t_1, t_2]$. Then for any $t, t'$ with $t_1 \le t < t' \le t_2$, we have $\tau(t) = \tau(t')$, hence
\[
\frac{\vp(t) + \ps(t)}{2} = \frac{\vp(t') + \ps(t')}{2} \implies (\vp(t') - \vp(t)) + (\ps(t') - \ps(t)) = 0
\]
Since $\vp$ and $\ps$ are non-decreasing, we conclude that $\vp(t) = \vp(t')$ and $\ps(t) = \ps(t')$. Thus both $\vp$ and $\ps$ are constant on $[t_1, t_2]$.


For any $s \in [0,1]$, define $\hvp(s) \coloneqq \vp(t)$ and $\hps(s) \coloneqq \psi(t)$ for any $t \in \tau^{-1}(s)$. Since $\vp$, $\psi$ are constant on the fibers $\tau^{-1}(s)$ of $\tau$, we obtain well-defined functions 
$\hvp \colon [0,1] \to [0,1]$ and $\hps \colon [0,1] \to [0,1]$ such that
\[
\vp(t) = \hvp(\tau(t)) \quad \text{and} \quad \ps(t) = \hps(\tau(t)) \quad \text{for all } t \in [0,1].
\]
By construction, $\g\circ\vp = \sigma\circ \psi$, so also $\g(\hvp(\tau(t))) = \s(\hps(\tau(t)))$ for all $t\in [0,1]$.
Surjectivity of $\tau$ then gives $\g\circ \hvp = \sigma\circ \hps$. Also, since $\varphi$, $\psi$ are surjective, so are $\hvp$ and $\hps$.

Since $\tau$ is a continuous map from a compact space ($[0,1]$) to a Hausdorff space (again $[0,1]$), it is a quotient map. By the universal property of quotient maps, since $\vp$ and $\ps$ are continuous and constant on the fibers of $\tau$, the induced maps $\hvp$ and $\hps$ are also continuous.

It remains to show strict monotonicity, which we do for $\hat\varphi$. Let $0 \le s_1 < s_2 \le 1$. We must show $\hvp(s_1) < \hvp(s_2)$.
Since $\tau$ is non-decreasing and surjective from $[0,1]$ to $[0,1]$, we can find $t_1, t_2 \in [0,1]$ such that $\tau(t_1) = s_1$, $\tau(t_2) = s_2$, and $t_1 < t_2$. Then
\[
\hat\varphi(s_1) = \varphi(t_1) \le \varphi(t_2) = \hat\varphi(s_2).
\]
Let us indirectly assume that equality holds: $\hvp(s_1) = \hvp(s_2)$. Then $\vp(t_1) = \vp(t_2)$ and since $\vp$ is non-decreasing, it must be constant on the entire interval $[t_1, t_2]$. By what was shown above, then $\ps$ must also be constant on $[t_1, t_2]$.
But then also $\tau(t)$ must be constant on $[t_1, t_2]$. In particular, $\tau(t_1) = \tau(t_2)$, which implies $s_1 = s_2$, a contradiction.
This proves that $\hvp$ is strictly monotonically increasing, as claimed.
\end{proof}

\end{document}
